% mage-auto-heading: the-mage-cycle
% Authoritative LaTeX manuscript; Markdown is a generated compatibility projection.

% <!-- part-foreshadows: modeling-principle, alignment-principle, govern-the-environment, failures-become-machinery -->

Chapters \magevalue{chapter:modeling|num} and \magevalue{chapter:alignment|num} separated Modeling from Alignment so that each mechanism could be examined clearly. Engineering does not usually present them that way. As work proceeds, new evidence changes what engineers understand about the problem; one representation exposes the need for another; implementation tests assumptions; and analysis changes the design. This chapter brings the two principles together as an engineering method: making consequential properties explicit, placing enforcement where obligations can be checked, and converting recurring judgment into durable engineering structure.

% <!-- principlebox -->

% <!-- box-family: canonical -->

\begin{quote}
\mbox{}%
\subsection{The MAGE cycle}\label{heading:the-mage-cycle}

Model the intent. Enforce stable obligations. Convert recurring judgment into durable engineering structure.

\textbf{Repeat.}
\end{quote}

The chapter follows that cycle from four directions. \magevalue{sec:the-dynamics-of-mage} develops the dynamics: models need not enter a system all at once; alignment can act first actively, then passively to preserve what it established; and because engineers often do not yet know the right models at the beginning of a project, implementation participates in discovering them. \magevalue{sec:interlude-one-problem-many-models} follows those dynamics through one problem in DocAble: a memory failure that drives successive changes in representation and execution, then a connected family of models, each answering a different question. \magevalue{sec:brownfield-engineering} enters a system that already exists, where the consequential structure arrived before the models that should describe it. \magevalue{sec:engineering-the-environment} turns to the destination. A task begins from the representations, controls, and evidence the environment already owns; work exposes what those structures can and cannot answer; and each useful conversion changes what later work inherits. Useful durable structure becomes engineering capital --- capital that can drift, depreciate, and cost more to maintain than it returns, so the chapter covers operation and retirement as well as accumulation.

The question here is practical: \textbf{what should you do next?}
